\chapter{我们的责任}\label{ch:g9:our-responsibility}

本书九年的主线在此汇合，归结为一个朴素的事实：
能读懂这一页的\omterm{def:g5:biodiversity:species}{物种}，正是唯一理解自身所属机制的\omterm{def:g5:biodiversity:species}{物种}——
而理解，
正如\cref{ch:g8:contraception}开篇所言，
带来选择。本章汇集理解所带来的种种责任：
对自身健康的责任，对他人健康的责任，
以及对承载我们所有人的\omterm{def:g6:exploring-our-environment:components}{生物}之网的
责任。

\section{对自身机制的责任}

\begin{proposition}[作为知情自我管理的健康]\label{prop:g9:our-responsibility:self}
本书关于健康的几条主线，汇总如下：
\begin{itemize}
  \item 日常契约
        （\cref{prop:g5:healthy-choices:contract}）——
        饮食、运动、\omterm{prop:g1:healthy-habits:needs}{睡眠}、清洁——是器官维护，
        如今可以逐一系统地说明其道理
        （\cref{prop:g7:heart-and-circulation:keep}、
        \cref{met:g7:lungs-and-gas-exchange:keep}、
        \cref{met:g8:protecting-the-nervous-system:manual}）；
  \item 各类物质的作用机制已经清楚
        （\cref{prop:g8:protecting-the-nervous-system:substances}），
        陷阱的不对称性也已权衡
        （\cref{rem:g5:healthy-choices:never}）；
  \item 生殖方面的选择，
        有了路线图和去处（\cref{ch:g8:contraception}）；
  \item 身体的\omterm{def:g9:heredity-and-traits:heredity}{遗传}史也已部分可知
        （\cref{ch:g9:heredity-and-traits}）：了解家族的倾向——
        早生的白发、易患的\omterm{def:g4:heart-and-blood:heart}{心脏}问题——并非命运，
        而是一份维护时间表
        （\cref{met:g9:unique-individuals:claims}：
        取值范围，需在生活中检验）。
\end{itemize}
\end{proposition}
\admitted

\begin{example}[从生物学看体检]\label{ex:g9:our-responsibility:checkup}
一次例行体检，就是本书内容在一次门诊中的重现：
测量\omterm{prop:g4:heart-and-blood:pulse}{脉搏}和血压（路边可读的循环网络，
\cref{ex:g7:heart-and-circulation:measures}），
反射锤（\cref{pb:g8:nervous-communication:1}），
疫苗接种记录（\cref{pb:g9:immune-defenses:1}），
家族病史（已部分可知的倾向），以及常规建议——
又是那份契约。医学在很大程度上，
是对你如今已掌握其说明书的机制所做的
应用性维护。
\end{example}

\section{对他人的责任}

\begin{proposition}[健康是共享的机制]\label{prop:g9:our-responsibility:others}
有三种机制使健康成为一笔共同的账：
\begin{itemize}
  \item \emph{传播}
        （\cref{def:g9:microbes-and-infection:stages}）：
        你的\omterm{def:g1:healthy-habits:hygiene}{卫生}习惯、咳嗽时的遮掩，以及患病时在家休息，
        都是对他人的保护；
  \item \emph{群体\omterm{def:g9:immune-defenses:memory}{免疫}}
        （\cref{prop:g9:immune-defenses:vaccination}）：
        \omterm{prop:g9:immune-defenses:vaccination}{接种疫苗}通过你，
        保护那些无法接种的人；
  \item \emph{医学的公共资源}：抗生素得以持续有效
        （\cref{exo:g9:how-species-change:10}），
        献血得以实施（\cref{ex:g9:unique-individuals:medicine}
        ——\omterm{prop:g4:journey-of-food:absorb}{血液}登记依赖献血者）——
        这些共享资源，其消耗或储备取决于个人的选择。
\end{itemize}
\end{proposition}
\admitted

\begin{example}[一位青少年的共同账目]\label{ex:g9:our-responsibility:account}
一周内的普通记录：发烧时待在家中，
切断了传播；加强针按时接种，维持了群体\omterm{def:g9:immune-defenses:memory}{免疫}；
春季完整服用一个疗程的抗生素，公共资源未被浪费；
十八岁时计划首次献血，充实了\omterm{prop:g4:journey-of-food:absorb}{血液}登记。
这些都算不上壮举——
这正是\cref{ex:g5:healthy-choices:day}的道理在一座城镇尺度上的体现：
公共\omterm{def:g1:healthy-habits:hygiene}{卫生}靠的是寻常日子里的坚持。
\end{example}

\section{对生物世界的责任}

\begin{proposition}[管理责任，附全部依据重述]\label{prop:g9:our-responsibility:stewardship}
自然主线及其机制汇总如下：
\begin{itemize}
  \item 每条\omterm{def:g3:food-chains:chain}{食物链}都始于绿色
        （\cref{prop:g6:plants-are-producers:firstlink}），
        每张餐桌都是某个生态系统的产出
        （\cref{prop:g6:producing-our-food:costs}）；
  \item \omterm{def:g5:food-webs:web}{食物网}借助多样性吸收冲击
        （\cref{prop:g5:biodiversity:web}），而多样性正是
        \omterm{prop:g9:how-species-change:selection}{自然选择}所依赖的储备
        （\cref{prop:g9:unique-individuals:reserve}、
        \cref{exo:g9:how-species-change:12}）——
        \omterm{def:g5:biodiversity:biodiversity}{生物多样性}不是风景，而是运营资本；
  \item \omterm{prop:g5:biodiversity:loss}{灭绝}是永久的
        （\cref{prop:g5:biodiversity:loss}）——
        唯一的有意例外天花，
        恰恰证明了这个词的绝对性
        （\cref{exo:g9:immune-defenses:12}）；
  \item 持久利用的规则已知
        （\cref{prop:g5:people-and-nature:lasting}），其违规
        行为可以诊断
        （\cref{ex:g7:oxygen-in-the-water:waste}），其修复方法
        也已演示
        （\cref{ex:g5:people-and-nature:river}、
        \cref{ex:g5:biodiversity:storks}）。
\end{itemize}
管理责任与健康一样，都是同一种维护思路，
只是作用于我们所生活其中的那个更大的身体。
\end{proposition}
\admitted

\begin{example}[公民的生物学有何用]\label{ex:g9:our-responsibility:citizen}
成年生活中的投票、购买与习惯，
总以本书所培养的问题形式不断出现：
这种渔业的捕捞是否快于资源的再生？这项
政策是保护\omterm{def:g5:flower-to-fruit:steps}{传粉}者的\omterm{def:g2:where-animals-live:habitat}{栖息地}，还是将其侵占？这一
所谓神奇饮食、疗法或“X基因”是否
自洽（\cref{met:g9:unique-individuals:claims}，
\cref{met:g8:contraception:evaluate}）？这种抗生素
是否应放入此饲槽？能进行这些审查的公民，
正是九年各章的用意所在。
\end{example}

\begin{method}[毕业生的方法]\label{met:g9:our-responsibility:method}
离开第 1 册时，应带走的是方法，而非仅仅是事实：
\begin{enumerate}
  \item 先观察，后下结论
        （\cref{met:g1:animals-around-us:observe} 开其先河，
        此后每次普查与笔记皆然）；
  \item 公平检验——只改变一个因素，设置对照
        （\cref{met:g4:what-plants-need:fairtest}，
        其后各年级皆沿用）；
  \item 指明机制——在其中的环节被指明之前，
        任何说法都不予审核
        （\cref{met:g8:contraception:evaluate}）；
  \item 权衡相互独立的佐证
        （\cref{met:g9:evidence-of-evolution:weigh}）；
  \item 并尊重生命的机制，无论在你的身体内
        还是身体之外：它古老、精妙，
        在一定限度内可以修复，也由你来守护。
\end{enumerate}
\end{method}

\begin{remark}[接下来的内容]\label{rem:g9:our-responsibility:next}
高中分册将以更高的放大率接续这些线索：
\omterm{def:g5:first-look-at-cells:cell}{细胞}的分子机器、附带完整数学的\omterm{def:g9:heredity-and-traits:heredity}{遗传}学、
深入展开的\omterm{def:g9:immune-defenses:memory}{免疫}系统与\omterm{def:g5:brain-in-command:system}{神经系统}、
作为定量科学的生态学，以及
作为其本然框架的\omterm{prop:g9:how-species-change:selection}{演化}。问题仍是
那些问题，只是追问得更加细致。带上方法。
\end{remark}

\section{练习}

\begin{exercise}[$\star$]\label{exo:g9:our-responsibility:1}
列出自我负责的四个方面，
见\cref{prop:g9:our-responsibility:self}。
\end{exercise}

\begin{exercise}[$\star$]\label{exo:g9:our-responsibility:2}
指出使健康得以共享的三种机制，
并各注明其所在的章。
\end{exercise}

\begin{exercise}[$\star$]\label{exo:g9:our-responsibility:3}
为何已知的家族倾向是一份计划表，而非
命运？
\end{exercise}

\begin{exercise}[$\star$]\label{exo:g9:our-responsibility:4}
结合其所在章节，阅读体检的三种手段。
\end{exercise}

\begin{exercise}[$\star$]\label{exo:g9:our-responsibility:5}
为何\omterm{def:g5:biodiversity:biodiversity}{生物多样性}是“营运资本”，而非
风景点缀？列出两种机制。
\end{exercise}

\begin{exercise}[$\star$]\label{exo:g9:our-responsibility:6}
凭记忆写出毕业生方法的五个步骤。
\end{exercise}

\begin{exercise}[$\star\star$]\label{exo:g9:our-responsibility:7}
论证在家休息的病假日属于一项公共\omterm{def:g1:healthy-habits:hygiene}{卫生}行为，
并指明其所属阶段。
\end{exercise}

\begin{exercise}[$\star\star$]\label{exo:g9:our-responsibility:8}
你的疫苗接种为谁筑起屏障？请以
群体作答，并列出两类受其保护的人。
\end{exercise}

\begin{exercise}[$\star\star$]\label{exo:g9:our-responsibility:9}
按\cref{ex:g9:our-responsibility:citizen}的方式审核：
“我们的新型喷雾剂能杀死果园中的每一种\omterm{prop:g3:sorting-living-things:groups}{昆虫}——
全面保护。”
\end{exercise}

\begin{exercise}[$\star\star$]\label{exo:g9:our-responsibility:10}
结合病房的选择，解释“抗生素是
一种公共资源”，并说明
维持其储备的两种习惯。
\end{exercise}

\begin{exercise}[$\star\star$]\label{exo:g9:our-responsibility:11}
将方法的五个步骤应用于一条\omterm{def:g9:microbes-and-infection:sorted}{病毒}式传播的“神奇
疗法”帖子，逐步进行。
\end{exercise}

\begin{exercise}[$\star\star\star$]\label{exo:g9:our-responsibility:12}
请你自己撰写本书的结语段落：
\omterm{def:g5:first-look-at-cells:cell}{细胞}、身体、\omterm{def:g5:food-webs:web}{食物网}、\omterm{def:g6:classification-kinship:tree}{亲缘树}各写一句，
并写出读者在这四者之中的位置。
\end{exercise}

\section{问题：毕业审核}

\begin{problem}[{周末问题——一个小镇的一年，
以全书的内容加以审核}]\label{pb:g9:our-responsibility:1}
你所在小镇（虚构）一年的新闻标题：冬季
流感流行；春季关于是否喷药防治悬铃木的争论；一
夏季公园池塘\omterm{prop:g3:sorting-living-things:groups}{鱼类}死亡；秋季新湿地保护区的
建设宣传活动；以及全年进行的
诊所疫苗接种与捐赠活动。

\medskip\noindent\textbf{第一部分——健康档案。}
\begin{enumerate}
  \item 流感流行：指出镇上建议所针对的传播途径，
        以及诊所为何
        优先为养老院员工\omterm{prop:g9:immune-defenses:vaccination}{接种疫苗}。
  \item 捐赠活动：说明献血登记册和疫苗接种台账
        各自记录了什么，以及
        谁会查阅这些记录。
  \item 一位专栏作家嘲笑“为洗手大惊小怪”。
        请以
        \cref{ex:g9:microbes-and-infection:history}中的病房为例，
        写一段文字作答。
  \item 药房反映，未服完的抗生素疗程又被退回。
        哪一种公共资源正被消耗，
        通过什么机制？
\end{enumerate}

\medskip\noindent\textbf{第二部分——自然档案。}
\begin{enumerate}[resume]
  \item 关于喷药的争论：运用果园审计
        （\cref{exo:g9:our-responsibility:9}）对
        “杀死所有\omterm{prop:g3:sorting-living-things:groups}{昆虫}”进行检验——还有谁会死亡，
        哪些免费服务会因此停止？
  \item \omterm{prop:g3:sorting-living-things:groups}{鱼类}死亡事件：炎热无风的一周，池塘
        呈豌豆绿色。请借助
        \cref{ch:g7:oxygen-in-the-water}进行诊断，
        并提出处理方案。
  \item 湿地保护运动：用三种机制为保护区辩护，
        即\omterm{def:g2:where-animals-live:habitat}{栖息地}、\omterm{def:g5:food-webs:web}{食物网}和多样性储备，
        每种各用一句话说明。
  \item 一位土地所有者问：“湿地对镇子究竟做过什么？”请以
        免费服务账目为依据作答
        （\cref{exo:g5:biodiversity:11}，
        以及河流对
        \cref{ex:g5:people-and-nature:river}中洪水的蓄滞作用）。
\end{enumerate}

\medskip\noindent\textbf{第三部分——毕业生的
结语。}
\begin{enumerate}[resume]
  \item 依据持续利用的规则，选出本年度最好和最差的决定，
        并论证每一项判断。
  \item 说明本方法的五个步骤如何在\omterm{prop:g3:sorting-living-things:groups}{鱼类}死亡事件的解决中发挥作用，
        即从最初的观察
        到被接受的机制。
  \item 以四个句子起草镇报的年终社论：
        健康共享、自然储备、
        主张经过审核、方法得以保留。
  \item 第 1 册最后一页的最后一行由你来写：
        用一句话写给
        从
        \cref{ch:g1:living-or-not}中那只蜗牛开始读起的读者。
\end{enumerate}
\end{problem}
